% Ad M. Brutum 1.13 --- headnote
% ad-brutum-01-13, 1 July 43 BC, in camp

\textit{M. Brutus to Cicero, written from his camp on
1 July 43 BC --- Perseus dateline \textit{Scr. in castris
K. Quint. a. 711 (43)}, confirmed by the closing
``\textit{Kal. Quintilibus ex castris}.'' The meta entry's
day-precision date is correct.}

\textit{The letter is a plea on behalf of children, written
in the moment when Brutus first suspects that his
brother-in-law M. Aemilius Lepidus is about to defect to
Antony. Lepidus's army in Narbonensis has been wavering
through May and June; the rumours have multiplied; Brutus
knows that if his sister Junia's marriage to Lepidus turns
into the marriage of a public enemy, his nephews will be
ruined by the senatorial decree of \textit{hostis
publicus}. He writes to Cicero, head of the senatorial
party, to put on the record now, before any vote falls,
that he speaks not as a brother-in-law but as the boys'
\textit{avunculus} (their maternal uncle), and that he
considers himself to have succeeded their father in
their care. The compression is unusual for Brutus, the
language unguarded (``my anxiety and my exasperation,''
\textit{sollicitudo ac stomachus}, are uncomfortably
candid). The fear was justified: the news of Lepidus's
junction with Antony reaches Rome on 30 June, the day
before this letter is written, and the Senate's vote
declaring him a public enemy will pass on 30 June or
shortly after, while this letter is on its way to Rome.}
